\documentclass[11pt]{amsart}
\usepackage{amsmath,amssymb}
\setlength{\emergencystretch}{3em}
\tolerance=600
\usepackage{tikz-cd}
\usepackage[hidelinks]{hyperref}

\newtheorem{thm}{Theorem}[section]
\newtheorem{lem}[thm]{Lemma}
\newtheorem{cor}[thm]{Corollary}
\newtheorem{prop}[thm]{Proposition}
\theoremstyle{definition}
\newtheorem{defi}[thm]{Definition}
\newtheorem{rem}[thm]{Remark}
\newtheorem{ex}[thm]{Example}
\newtheorem{que}[thm]{Question}

\newcommand{\Same}{\mathrm{Same}}
\newcommand{\Presh}{\mathrm{Presh}}
\newcommand{\Hom}{\mathrm{Hom}}
\newcommand{\im}{\mathrm{im}}
\newcommand{\Sh}{\mathrm{Sh}}
\newcommand{\Sub}{\mathrm{Sub}}
\newcommand{\Set}{\mathbf{Set}}
\newcommand{\cC}{\mathcal C}
\newcommand{\cE}{\mathcal E}
\newcommand{\Nuc}{\mathrm{Nuc}}

\begin{document}

\title[Joins of localisations and the sameness lattice]{Joins of localisations in the lattice of two-out-of-three classes: the finite-iteration criterion for closure operators on exact categories, with applications to toposes, torsion theories and fuzzy sets}
\author{Ryota Shibusawa}
\thanks{Corresponding author. E-mail: \texttt{r-shibusawa@daiichi-koudai.ac.jp}}
\address{Daiichi Institute of Technology}
\email{r-shibusawa@daiichi-koudai.ac.jp}
\keywords{closure operator, exact category, Lawvere--Tierney topology, hereditary torsion theory, Gabriel localisation, nucleus, sheafification, two-out-of-three, fuzzy set, join of localisations, finite iteration}
\subjclass[2020]{18E08 (primary), 18F10, 18E40, 18B25, 06D22, 03E72 (secondary)}

\begin{abstract}
A closure operator $c$ on the subobjects of an exact category, stable under
pullback and preserving finite intersections, determines the class $W_c$ of
\emph{$c$-bidense} maps, those whose image is $c$-dense and whose kernel
pair lies in the closure of the diagonal. For a Lawvere--Tierney topology on
a topos this is the class of maps inverted by sheafification; for a
hereditary torsion theory on a Grothendieck abelian category, whose
closure operators are exactly the hereditary torsion theories, it is the
class of maps with torsion kernel and cokernel, inverted by the Gabriel
quotient; for a nucleus on a frame viewed as a category it is the class of
inequalities identified by the nucleus. These classes are elements of the
complete lattice $\Same(\cC)$ of classes of morphisms closed under
composition and two-out-of-three. We prove that for two closure operators
$c,c'$ the join $W_c\vee W_{c'}$ is the class of maps whose image and kernel
pair are reached by a bounded number of alternating applications of $c$ and
$c'$, and that every such map is a composite of maps of $W_c\cup W_{c'}$, so
that the join is the closure of the union under composition alone. In a
Grothendieck topos and in a Grothendieck abelian category the least closure
operator $\bar c$ above $c$ and $c'$ is the transfinite iterate of $cc'$, and
$W_c\vee W_{c'}=W_{\bar c}$ if and only if $(cc')^k=\bar c$ for some $k$: for
subtoposes this is the identity $(jj')^k=j\vee j'$ on the subobject
classifier, for torsion theories $T,T'$ it says that every object of the
join torsion class has an alternating $T'/T$-filtration of bounded length.
For persistence modules, representations of the ordinal $\omega$, with the
torsion classes of modules supported on the even and on the odd vertices,
the join torsion class is everything but the interval modules need
filtrations of unbounded length, so the join of the classes of maps that
are isomorphisms at the odd vertices and at the even vertices is a proper
class, the linear analogue of the known failure for sheaves on $\mathbb N$.
In a frame viewed as a poset the witnesses that force uniformity are
absent and the join is governed by the pointwise identity, as it is for
isomorphisms up to a nucleus of fuzzy sets by a monomorphism version of
the same argument; the paper thus places the joins of sheaf-equivalence
classes, of torsion-theoretic localisations, of nuclei on a frame and of
isomorphisms up to a nucleus of fuzzy sets under one theorem.
\end{abstract}

\maketitle

\section{Introduction}

Let $\cC$ be a category. A \emph{sameness} on $\cC$ is a class of morphisms
containing the identities and closed under composition and two-out-of-three;
samenesses form a complete lattice $\Same(\cC)$ under inclusion, meets being
intersections and the join of two classes the least sameness containing both
\cite{Same}. Every localisation of $\cC$, in particular every reflective
subcategory, contributes the class $W$ of maps it inverts. For a family of
localisations closed under joins, the question studied in
\cite{Topjoin,Finitejoin,Stabjoin} is whether $W$ sends the join of two
localisations to the join of the two classes: whether every map inverted by
the join localisation is obtained from maps inverted by the two
localisations by finitely many compositions and cancellations. For subtoposes
of a presheaf topos the answer was known to be yes on every finite category
\cite{Finitejoin} and no on the poset $\mathbb N$ \cite{Topjoin}, and
\cite{Stabjoin} gave the exact criterion for subtoposes of a Grothendieck
topos: $W_j\vee W_{j'}=W_{j\vee j'}$ iff $(jj')^k=j\vee j'$ for some $k$, an
identity between endomorphisms of the subobject classifier. The same
identity, $(j_2j_1)^N=j_1\vee j_2$ for nuclei on a Heyting algebra, is the
finite-iteration axiom of the lax logic of the join of two nuclei
\cite{Laxjoin}, and it is the exact condition for the join of the classes
of isomorphisms up to the two nuclei in the sameness lattice of fuzzy sets
to be the class of the join nucleus \cite{Fuzzysame}, whereas in the
sameness lattice of the Heyting algebra itself the pointwise identity
suffices.

This paper proves one theorem behind these results and applies it to a
setting not covered by them, the Gabriel localisations of abelian
categories. The common structure is a \emph{closure operator} $c$ on the
subobjects of an exact category $\cC$: a pullback-stable,
intersection-preserving closure on each $\Sub(X)$. It determines the
sameness $W_c$ of \emph{$c$-bidense} maps, the maps $f\colon X\to Y$ with
$c(\im f)=Y$ and $\ker f\subseteq c(\Delta_X)$. For two closure operators
$c,c'$ let $U_k$ be the class of maps with
\[
(cc')^k(\im f)=Y\qquad\text{and}\qquad\ker f\subseteq(cc')^k(\Delta_X),
\]
the maps that are bidense \emph{uniformly} for the $k$-fold alternation of
the two closures.

\begin{thm}\label{thm:A}
Let $c,c'$ be closure operators on an exact category $\cC$. Then
\[
W_c\vee W_{c'}=\bigcup_kU_k
\]
in $\Same(\cC)$; $U_0$ is the class of isomorphisms, and for $k\ge1$ every
map in $U_k$ is a composite of $2k$ regular epimorphisms followed by $2k$
monomorphisms lying alternately in $W_{c'}$ and $W_c$. In particular the
join is the closure of $W_c\cup W_{c'}$ under composition.
\end{thm}

\begin{thm}\label{thm:B}
Let $\cC$ be a Grothendieck topos or a Grothendieck abelian category, and
$c,c'$ closure operators on $\cC$. Then the least closure operator $\bar c$
above $c$ and $c'$ is the transfinite iterate of $cc'$, and
$W_c\vee W_{c'}=W_{\bar c}$ if and only if $(cc')^k=\bar c$ for some $k$.
\end{thm}

Theorem~\ref{thm:B} is proved under axioms (Section~\ref{sec:criterion})
that isolate what is used: chains of subobjects have pullback-stable
unions, and the \emph{fold map} of two copies of an object glued along a
subobject $S$ witnesses whether $S$ is reached in $k$ steps; the two classes
of categories named in the statement are the ones in which we verify the
axioms, by different arguments.

In a Grothendieck topos the closure operators are the Lawvere--Tierney
topologies and $W_c$ is the class of maps inverted by sheafification; this is
the main theorem of \cite{Stabjoin}, whose proofs are reproduced here in the
generality in which they were written, and whose consequences for
presheaf toposes, ideal topologies, locales and subspaces we summarise
(Section~\ref{sub:topos}, \ref{sub:locale}). In a Grothendieck abelian
category $\mathcal A$ the closure operators are exactly the hereditary
torsion theories $T$ (Proposition~\ref{prop:torsion}): $c_T(A)$ is the
preimage of the torsion part of $X/A$, $W_{c_T}$ is the class of maps with
kernel and cokernel in $T$, the maps inverted by the quotient functor
$\mathcal A\to\mathcal A/T$, and $\bar c$ is the closure operator of the join
torsion theory $T\vee T'$. Theorems~\ref{thm:A} and~\ref{thm:B} read:

\begin{cor}\label{cor:abelian}
Let $T,T'$ be hereditary torsion theories on a Grothendieck abelian category
$\mathcal A$ and let $E_k$ be the class of objects having a filtration
$0=M_0\subseteq M_1\subseteq\cdots\subseteq M_{2k}=M$ with $M_{2i+1}/M_{2i}
\in T'$ and $M_{2i+2}/M_{2i+1}\in T$. Then $W_T\vee W_{T'}$ is the class of
maps whose kernel and cokernel lie in $\bigcup_kE_k$, the closure of
$W_T\cup W_{T'}$ under composition; and $W_T\vee W_{T'}=W_{T\vee T'}$ iff
$T\vee T'=E_k$ for some $k$, i.e.\ iff the objects of the join torsion class
have alternating $T'/T$-filtrations of bounded length.
\end{cor}

For persistence modules, i.e.\ representations $\cdots\to M_2\to M_1\to M_0$
of the ordinal $\omega$ in vector spaces, with $T$ and $T'$ the modules
supported on the even and on the odd vertices, $W_T$ is the class of maps
that are isomorphisms at the odd vertices, $T\vee T'$ is everything, and the
interval module $k[0,N]$ needs an alternating filtration of length $N+2$;
so $W_T\vee W_{T'}$ is a proper class, the class of maps whose kernel and
cokernel are ``bounded in alternation'' (Example~\ref{ex:persistence}). This
is the linear analogue of the failure for sheaves on $\mathbb N$
\cite{Topjoin}, and, as far as we know, the first statement of this kind
about Gabriel localisations.

The frame $L$ itself, viewed as a poset category, is an exact category in
which every morphism is a monomorphism, with closure operators
$c_n(x\le y)=n(x)\wedge y$ for nuclei $n$; Theorem~\ref{thm:A} applies and
gives $W_n\vee W_{n'}=\{x\le y:(nn')^kx\ge y\text{ for some }k\}$, so that the
join is the class of the join nucleus iff the identity $(nn')^kx=(n\vee n')x$
holds \emph{pointwise}; the witnesses of Theorem~\ref{thm:B} do not exist in
a poset, and indeed the uniform identity is not needed there
(Section~\ref{sub:frame}). Fuzzy sets are handled by a monomorphism version
of Theorem~\ref{thm:A} (Proposition~\ref{prop:mono}), since the maps
inverted by the reflection onto $j$-closed fuzzy sets are bijections; this
recovers the description of their joins by bounded alternating chains
\cite{Fuzzysame} (Section~\ref{sub:fuzzy}).

The proofs are elementary and use only the calculus of closure operators on
subobjects. The inclusion $U_k\subseteq W_c\vee W_{c'}$ is a peeling argument:
a regular epimorphism whose kernel pair lies in the $k$-fold alternating
closure of the diagonal is divided by the equivalence relation ``kernel pair
intersected with the $c'$-closure of the diagonal'', which is a
$c'$-bidense map, and the induced map needs one closure fewer; dually a
monomorphism is enlarged to the $c'$-closure of its image. The failure
direction of Theorem~\ref{thm:B} uses, for a $\bar c$-dense subobject $S$
of $Y$ that is not reached in $k$ steps, the fold map of two copies of $Y$
glued along $S$; the coproduct of these maps over all $k$ is $\bar c$-bidense
but lies in no $U_k$. The criterion is uniform in the subobject, and it is
the iteration of the closure operators, not of the reflectors, that
matters: on a monoid of order $8$ the alternating sheafification tower
never stabilises although the join is preserved
(Remark~\ref{rem:tower}).

The quotient-closed regime is different. For subvarieties $R,R'$ of a
variety of algebras, $W_R\vee W_{R'}=W_{R\cap R'}$ always holds, because the
composite of the two reflectors is the reflector onto $R\cap R'$, while the
meet $W_R\cap W_{R'}$ is the class of the join variety only under a fracture
condition \cite{Varieties}; the natural closure operators there are not
pullback-stable, and we explain in Section~\ref{sec:dual} why no common
mechanism exists.

\subsection*{What is known.}
Samenesses and the join theorem for a composite reflector are from
\cite{Same,Postnikov}; the meet theorem, the reduction lemma and the
$\mathbb N$ counterexample from \cite{Topjoin}; the finite case from
\cite{Finitejoin}; the criterion for subtoposes of Grothendieck toposes,
with the presheaf, ideal-type, object-generated, localic and subspace
cases, from \cite{Stabjoin}. Closure operators on categories are the subject
of \cite{DT}; universal closure operators on a topos, their correspondence
with Lawvere--Tierney topologies and the bidense description of the maps
inverted by sheafification are in \cite[V.1, V.3]{MLM},
\cite[A4.3--A4.4]{Elephant}. Hereditary torsion theories, Gabriel's quotient
categories and the characterisation of the maps they invert are in
\cite{Gabriel,Stenstrom}; the correspondence between hereditary torsion
theories and closure operators on module categories is discussed in
\cite{DT}, and its extension to normal and homological categories in
\cite{CDT}; the join of hereditary torsion theories and its description
by transfinite extensions go back to \cite{Dickson}. Joins of nuclei as transfinite iterates are in \cite{Escardo};
the weak topologies of \cite{KM} compute the join of two Lawvere--Tierney
topologies as such an iterate. The lax logic of the join of two nuclei and
its finite-iteration axiom are in \cite{Laxjoin}; the sameness lattice of
fuzzy sets and the uniform and pointwise criteria in \cite{Fuzzysame}; the
quotient-closed regime in \cite{Varieties}.

\subsection*{What is new.}
The main new result is Corollary~\ref{cor:abelian}: for hereditary torsion
theories $T,T'$ on a Grothendieck abelian category, the join of the classes
of maps inverted by the two Gabriel quotients is the class of maps whose
kernel and cokernel have finite alternating $T'/T$-filtrations, it is
generated by composition alone, and it is the class of the join torsion
theory exactly when the lengths are bounded; Example~\ref{ex:persistence}
shows that the bound can fail even when the join torsion theory is the
whole category. The description of $T\vee T'$ by transfinite alternating
filtrations and the dictionary between closure operators and hereditary
torsion theories (Proposition~\ref{prop:torsion}) are classical; the
statement about the join of the inverted classes, and its failure, are
new as far as we know. The general Theorems~\ref{thm:A} and~\ref{thm:B}, in
exact categories with closure operators and with the witness axioms
verified in toposes and in abelian categories, are the tool; the
topos-theoretic, localic, frame and fuzzy-set statements of
\cite{Stabjoin,Fuzzysame} are recovered as instances, with the
pointwise/uniform difference explained by the presence or absence of
witnesses (Section~\ref{sec:instances}), and the quotient-closed regime is
contrasted in Section~\ref{sec:dual}.

\section{Closure operators on exact categories}\label{sec:prelim}

Throughout, $\cC$ is an \emph{exact} category in the sense of Barr
\cite{Barr}: it has finite limits, every morphism factors as a regular
epimorphism followed by a monomorphism, regular epimorphisms are stable under
pullback, and every equivalence relation is the kernel pair of its
coequaliser. We write $\im f\subseteq Y$ for the image of $f\colon X\to Y$,
$\ker f\subseteq X\times X$ for its kernel pair, $\Delta_X\subseteq X\times X$
for the diagonal, $g^{-1}(B)$ for the pullback of a subobject and $g(A)$ for
the image of a subobject. We use freely the calculus of subobjects in a
regular category: $\im(gh)=g(\im h)$; $\ker(gh)=(h\times h)^{-1}(\ker g)$;
for a regular epimorphism $e$, $e(e^{-1}(B))=B$; products of regular
epimorphisms are regular epimorphisms; and for $g\colon Y\to Z$ and
$A\subseteq Y$,
\[
g^{-1}(g(A))=\pi_1\bigl(\ker g\cap(Y\times A)\bigr),
\]
as both are the image of $Y\times_ZA\to Y$.

\begin{defi}
A \emph{closure operator} on $\cC$ assigns to each object $X$ a map
$A\mapsto c(A)$ on $\Sub(X)$ that is inflationary, monotone, idempotent,
preserves binary intersections, and is stable under pullback:
$g^{-1}(c(B))=c(g^{-1}(B))$ for every $g\colon X\to Y$ and $B\subseteq Y$.
Closure operators are ordered pointwise. A subobject $A\subseteq X$ is
\emph{$c$-dense} if $c(A)=X$.
\end{defi}

These are idempotent closure operators in the sense of \cite{DT}, with
the additional requirements that they preserve binary intersections and
commute with pullbacks along all morphisms (the hereditary property of
\cite{DT} concerns pullbacks along monomorphisms only). On a topos they are
exactly the universal closure operators, which correspond to the
Lawvere--Tierney topologies \cite[V.1]{MLM}; on a Grothendieck abelian
category they are exactly the hereditary torsion theories
(Proposition~\ref{prop:torsion}). A \emph{word} is a composite
$w=c_1c_2\cdots c_n$ of closure operators, $c_n$ applied first; words are
inflationary, monotone, pullback-stable and preserve binary intersections,
but need not be idempotent.

\begin{lem}\label{lem:calc}
Let $w$ be a word, $g\colon X\to Y$, $A,K\subseteq X$, $B\subseteq Y$.
\begin{enumerate}
\item[(i)] $g(w(A))\subseteq w(g(A))$.
\item[(ii)] $w(A)\cap K\subseteq w(A\cap K)$.
\item[(iii)] $w(A\times B)=w(A)\times w(B)$ in $\Sub(X\times Y)$.
\item[(iv)] For a closure operator $c$, $c(\Delta_X)$ is an equivalence
relation on $X$.
\end{enumerate}
\end{lem}
\begin{proof}
It suffices to treat a single closure operator $c$ in (i)--(iii) and to
iterate. (i) $A\subseteq g^{-1}(g(A))$, so $c(A)\subseteq c(g^{-1}(g(A)))=
g^{-1}(c(g(A)))$. (ii) Let $m\colon K\hookrightarrow X$; then $A\cap K=
m(m^{-1}(A))$ and $c(A)\cap K=m(m^{-1}(c(A)))=m(c_K(m^{-1}(A)))\subseteq
c(m(m^{-1}(A)))$ by (i), where $c_K$ is the closure on $\Sub(K)$, the
pullback of $c$. (iii) $A\times B=\pi_1^{-1}(A)\cap\pi_2^{-1}(B)$, and $c$
preserves intersections and pullbacks. (iv) Reflexivity and symmetry are
clear (the twist map preserves $\Delta_X$ and commutes with $c$). The
composite of the relation $c(\Delta_X)$ with itself is
$\pi_{13}\bigl(\pi_{12}^{-1}c(\Delta_X)\cap\pi_{23}^{-1}c(\Delta_X)\bigr)=
\pi_{13}\bigl(c(\pi_{12}^{-1}\Delta_X\cap\pi_{23}^{-1}\Delta_X)\bigr)\subseteq
c(\pi_{13}(\delta))=c(\Delta_X)$ by (i), where $\delta\subseteq X^3$ is the
image of the triple diagonal.
\end{proof}

\begin{defi}
A map $f\colon X\to Y$ is \emph{$c$-bidense} if $\im f$ is $c$-dense in $Y$
and $\ker f\subseteq c(\Delta_X)$. Write $W_c$ for the class of $c$-bidense
maps.
\end{defi}

A monomorphism is $c$-bidense iff it is $c$-dense, and a regular
epimorphism iff its kernel pair lies in $c(\Delta_X)$. On a topos with the
closure operator of a Lawvere--Tierney topology $j$, $W_c$ is the class of
maps inverted by the sheafification $a_j$ \cite[A4.3]{Elephant}; in general
$W_c$ is defined by the closure operator alone, and Lemma~\ref{lem:U}
below (with $c'=c$) shows that it is a sameness.

\section{Uniform bidensity and the join}\label{sec:uniform}

Fix closure operators $c,c'$ on $\cC$ and write $c^{(k)}:=(cc')^k$, a word
of length $2k$ ending with $c'$, $c^{(0)}$ being the identity; since
closures are inflationary, $c^{(k)}\le c^{(k+1)}$ and $c,c'\le c^{(1)}$.
Write $V:=W_c\vee W_{c'}$.

\begin{defi}
For $k\ge0$ let $U_k$ be the class of maps $f\colon X\to Y$ with
$c^{(k)}(\im f)=Y$ and $\ker f\subseteq c^{(k)}(\Delta_X)$. Put
$U=\bigcup_kU_k$.
\end{defi}

$U_0$ is the class of isomorphisms: $\im f=Y$ means that the monic part of
$f$ is an isomorphism, and $\ker f=\Delta_X$ that $f$ is a monomorphism, so
$f$ is a monic regular epimorphism.

\begin{lem}\label{lem:U}
$U_k\subseteq U_{k+1}$, $W_c\cup W_{c'}\subseteq U_1$, and $U$ is a sameness.
Hence $V\subseteq U$. With $c'=c$ the same proof shows that $W_c$ is a
sameness.
\end{lem}
\begin{proof}
The first two assertions follow from $c^{(k)}\le c^{(k+1)}$ and
$c,c'\le c^{(1)}$. Identities lie in $U_0$. Let $h\colon X\to Y$ and
$g\colon Y\to Z$; we use Lemma~\ref{lem:calc} for the words $c^{(k)}$.

\emph{Composition.} Let $h\in U_k$, $g\in U_l$. Then $\im(gh)=g(\im h)$ and
$c^{(k+l)}(\im gh)\supseteq c^{(l)}\bigl(g(c^{(k)}(\im h))\bigr)=c^{(l)}(g(Y))
=c^{(l)}(\im g)=Z$ by (i). And $\ker(gh)=(h\times h)^{-1}(\ker g)\subseteq
(h\times h)^{-1}(c^{(l)}(\Delta_Y))=c^{(l)}((h\times h)^{-1}\Delta_Y)=
c^{(l)}(\ker h)\subseteq c^{(l+k)}(\Delta_X)$.

\emph{Cancellation of $g$.} Let $gh\in U_k$ and $g\in U_k$. Then
$\ker h\subseteq\ker(gh)\subseteq c^{(k)}(\Delta_X)$. Let $B:=g^{-1}(\im gh)
\subseteq Y$; then $c^{(k)}(B)=g^{-1}(c^{(k)}(\im gh))=g^{-1}(Z)=Y$. Moreover
$B=\pi_1\bigl(\ker g\cap(Y\times\im h)\bigr)$, and
\begin{align*}
\pi_1\bigl(\ker g\cap(Y\times\im h)\bigr)&\subseteq
\pi_1\bigl(c^{(k)}(\Delta_Y)\cap(Y\times\im h)\bigr)\\
&\subseteq\pi_1\bigl(c^{(k)}(\Delta_Y\cap(Y\times\im h))\bigr)\\
&\subseteq c^{(k)}\bigl(\pi_1(\Delta_Y\cap(Y\times\im h))\bigr)=c^{(k)}(\im h)
\end{align*}
by (ii) and (i). Hence $Y=c^{(k)}(B)\subseteq c^{(2k)}(\im h)$.

\emph{Cancellation of $h$.} Let $gh\in U_k$ and $h\in U_k$. Then
$\im g\supseteq\im(gh)$, so $c^{(k)}(\im g)=Z$. Since $h\times h\colon
X\times X\to\im h\times\im h$ is a regular epimorphism,
$\ker g\cap(\im h\times\im h)=(h\times h)(\ker gh)\subseteq
(h\times h)(c^{(k)}(\Delta_X))\subseteq c^{(k)}(\Delta_Y)$ by (i); and
$c^{(k)}(\im h\times\im h)=Y\times Y$ by (iii). So
$\ker g=\ker g\cap c^{(k)}(\im h\times\im h)\subseteq
c^{(k)}\bigl(\ker g\cap(\im h\times\im h)\bigr)\subseteq c^{(2k)}(\Delta_Y)$
by (ii).

For $c'=c$ one has $c^{(k)}=c$ for all $k\ge1$, so $U_k=W_c$ and the
argument shows that $W_c$ is closed under composition and cancellation.
\end{proof}

\begin{lem}[peeling]\label{lem:peel}
Let $w=c_1c_2\cdots c_n$ be a word in closure operators.
\begin{enumerate}
\item[(a)] If $f\colon X\to Y$ is a regular epimorphism with $\ker f
\subseteq w(\Delta_X)$, then $f=f_1\circ\cdots\circ f_n\circ u$ with $u$ an
isomorphism and $f_i$ a regular epimorphism in $W_{c_i}$.
\item[(b)] If $f\colon X\to Y$ is a monomorphism with $w(X)=Y$, then
$f=f_1\circ\cdots\circ f_n\circ u$ with $u$ an isomorphism and $f_i$ a
monomorphism in $W_{c_i}$.
\end{enumerate}
For $n\ge1$ the isomorphism $u$ can be absorbed into $f_n$.
\end{lem}
\begin{proof}
Induction on $n$; for $n=0$ the hypothesis says that $f$ is an isomorphism.
Write $w=w'c_n$.
(a) Let $R:=\ker f\cap c_n(\Delta_X)$, an equivalence relation on $X$ by
Lemma~\ref{lem:calc}(iv), and let $q\colon X\twoheadrightarrow X/R$ be its
coequaliser, which has $\ker q=R$ by exactness; so $\ker q\subseteq
c_n(\Delta_X)$ and $q\in W_{c_n}$. Since $R\subseteq\ker f$, $f$ coequalises
the two projections of $R$ and $f=\bar f\circ q$ with $\bar f$ a regular
epimorphism, and $\ker\bar f=(q\times q)(\ker f)$ because $q\times q$ is a
regular epimorphism and $\ker f=(q\times q)^{-1}(\ker\bar f)$. Now
\begin{align*}
(q\times q)(\ker f)&=(q\times q)\bigl(w'(c_n\Delta_X)\cap\ker f\bigr)
\subseteq(q\times q)\bigl(w'(c_n\Delta_X\cap\ker f)\bigr)\\
&=(q\times q)(w'(R))\subseteq w'\bigl((q\times q)(R)\bigr)=w'(\Delta_{X/R})
\end{align*}
by Lemma~\ref{lem:calc}(ii) and (i), as $(q\times q)(\ker q)=\Delta_{X/R}$.
Induction applied to $\bar f$ and $w'$ gives the factorisation.
(b) Regard $X$ as a subobject of $Y$ and let $Z:=c_n(X)\subseteq Y$. The
inclusion $X\hookrightarrow Z$ is $c_n$-dense (the closure of $X$ in $Z$ is
$c_n(X)\cap Z=Z$ by pullback stability), hence in $W_{c_n}$, and
$w'(Z)=w(X)=Y$. Induction applied to $Z\hookrightarrow Y$ and $w'$ gives the
factorisation.
\end{proof}

\begin{proof}[Proof of Theorem~\ref{thm:A}]
Let $C$ be the closure of $W_c\cup W_{c'}$ under composition. Then
$C\subseteq V\subseteq U$ by Lemma~\ref{lem:U}. Let $f\in U_k$, $k\ge1$, and
factor $f=m\circ e$ with $e\colon X\twoheadrightarrow\im f$ a regular
epimorphism and $m\colon\im f\hookrightarrow Y$ a monomorphism. Then
$\ker e=\ker f\subseteq c^{(k)}(\Delta_X)$ and $c^{(k)}(\im f)=Y$, and
$c^{(k)}$ is a word of length $2k$ ending with $c'$. By
Lemma~\ref{lem:peel}, $e$ is a composite of $2k$ regular epimorphisms and
$m$ a composite of $2k$ monomorphisms, alternately in $W_{c'}$ and $W_c$,
beginning with $W_{c'}$. So $U\subseteq C$, and $C=V=U$.
\end{proof}

\begin{figure}[ht]
\[
\begin{tikzcd}[column sep=normal,row sep=large]
X\arrow[r,two heads,"W_{c'}"]\arrow[rrr,bend left=25,"e"] & X/R_1\arrow[r,two heads,"W_c"] & \cdots\arrow[r,two heads,"W_c"] & \im f\arrow[d,hook,"W_{c'}"]\arrow[ddlll,bend left=35,"m"{near start}]\\
& & & Z_1\arrow[d,hook,"W_c"]\\
Y & \cdots\arrow[l,hook',"W_c"'] & & Z_2\arrow[ll,hook',"W_{c'}"']
\end{tikzcd}
\]
\caption{Theorem~\ref{thm:A}: a map in $U_k$ is $2k$ quotients by ``kernel
pair intersected with the closure of the diagonal'' followed by $2k$
alternating closures of its image, each step bidense for one of the two
closure operators.}
\end{figure}

\begin{rem}\label{rem:cancel}
In general the join of two samenesses requires the two-out-of-three
saturation of the compositional closure \cite{Same}: on the poset $0<1<2$
viewed as a category, the classes $W_1=\{0\to1\}$ and $W_2=\{0\to2\}$ (with
the identities) are samenesses, their compositional closure is
$\{0\to1,0\to2\}$, and two-out-of-three forces $1\to2$ into the join, which
is therefore the class of all morphisms. For two bidense classes no
cancellation is needed. The number of steps is controlled in both
directions: a composite of $n$ maps of $W_c\cup W_{c'}$ lies in $U_n$, and a
map of $U_k$ is a composite of $4k$ such maps.
\end{rem}

\begin{prop}[monomorphisms]\label{prop:mono}
Let $\cC$ be a category with finite limits and $c,c'$ closure operators on
$\cC$ (pullback-stable, intersection-preserving closures on the subobject
lattices); exactness is not assumed. For monomorphisms $m\colon X\hookrightarrow
Y$ let $M_k$ be the class of those with $c^{(k)}(X)=Y$ and $M=\bigcup_kM_k$.
Then $M$ contains the $c$-dense and the $c'$-dense monomorphisms, is closed
under composition, and under cancellation among monomorphisms: if $gh\in M_k$
then $h\in M_k$ and $g\in M_k$; and every $m\in M_k$ is
a composite of $2k$ monomorphisms alternately $c'$-dense and $c$-dense.
Consequently, if $B$ is a sameness consisting of monomorphisms such that for
$m\colon X\hookrightarrow Y$ in $B$ the inclusions $X\hookrightarrow c(X)$ and
$X\hookrightarrow c'(X)$ are in $B$, then the join in $\Same(\cC)$ of the
$c$-dense and the $c'$-dense members of $B$ is $M\cap B$.
\end{prop}
\begin{proof}
For $h\colon X\hookrightarrow Y$ in $M_k$ and $g\colon Y\hookrightarrow Z$ in
$M_l$, pullback stability along $Y\hookrightarrow Z$ gives
$c^{(k)}_Z(X)\cap Y=c^{(k)}_Y(X)=Y$, so $c^{(k)}_Z(X)\supseteq Y$ and
$c^{(k+l)}_Z(X)\supseteq c^{(l)}_Z(Y)=Z$. If $gh\in M_k$, then $c^{(k)}_Y(X)=c^{(k)}_Z(X)\cap Y=Z\cap Y=Y$, so
$h\in M_k$, and $c^{(k)}_Z(Y)\supseteq c^{(k)}_Z(X)=Z$, so $g\in M_k$; neither
cancellation needs the other factor to be in $M$. The
factorisation is Lemma~\ref{lem:peel}(b), whose proof uses only pullback
stability. For the last assertion, $M\cap B$ is a sameness containing the
two classes, and each $m\in M_k\cap B$ is a composite of its peeling
factors, which lie in $B$ by hypothesis and by two-out-of-three in $B$.
\end{proof}

\section{The finite-iteration criterion}\label{sec:criterion}

In this section $\cC$ satisfies
\begin{enumerate}
\item[(H1)] every $\Sub(X)$ is a set and has unions of chains, these are
stable under pullback, and $A\cap\bigcup_\alpha B_\alpha=\bigcup_\alpha(A\cap
B_\alpha)$ for chains $(B_\alpha)$.
\end{enumerate}
Let $c,c'$ be closure operators and define the transfinite iterates
$c^{(0)}(A)=A$, $c^{(\alpha+1)}(A)=c^{(1)}(c^{(\alpha)}(A))$ and
$c^{(\lambda)}(A)=\bigcup_{\alpha<\lambda}c^{(\alpha)}(A)$ for limit
$\lambda$; they increase with $\alpha$, are constant from some ordinal
$\gamma_A$ on since $\Sub(X)$ is a set, commute with pullback by (H1), and
$c^{(\alpha)}(c^{(\beta)}(A))=c^{(\beta+\alpha)}(A)$.

\begin{lem}\label{lem:joinK}
Under (H1), $\bar c(A):=\bigcup_\alpha c^{(\alpha)}(A)$ is a closure
operator, the least one above $c$ and $c'$.
\end{lem}
\begin{proof}
$\bar c$ is inflationary and monotone; $\bar c(A)=c^{(\gamma)}(A)$ for
$\gamma=\gamma_A$ is a fixed point of $c^{(1)}$, hence of $c$ and $c'$ (as
$c'(B)\subseteq c(c'(B))=B$ for a fixed point $B$) and of $\bar c$; it
commutes with pullback because each $c^{(\alpha)}$ does and pullback
preserves unions of chains; and it preserves binary intersections because
each $c^{(\alpha)}$ does and, by (H1),
$\bigcup_\alpha c^{(\alpha)}(A)\cap\bigcup_\beta c^{(\beta)}(B)=
\bigcup_{\alpha,\beta}c^{(\alpha)}(A)\cap c^{(\beta)}(B)=\bigcup_\gamma
c^{(\gamma)}(A\cap B)$. It lies above $c$ and $c'$, and if $d$ is a closure
operator above $c$ and $c'$ then $c^{(\alpha)}(A)\subseteq d(A)$ for all
$\alpha$ by induction, $d$ being monotone and idempotent, with unions at
limits.
\end{proof}

A map in $U_k$ is $\bar c$-bidense, since $c^{(k)}\le\bar c$; so
$V\subseteq U\subseteq W_{\bar c}$. The converse inclusion $W_{\bar c}\subseteq U$
is what the criterion is about, and its failure is detected by
\emph{witnesses}.

\begin{defi}\label{def:witness}
A \emph{witness} for a subobject $S\subseteq Y$ is a regular epimorphism
$f_S\colon X_S\to Y$ such that, for every word $w$ and for $w=\bar c$,
$\ker f_S\subseteq w(\Delta_{X_S})$ iff $w(S)=Y$. A \emph{sum} of a sequence
of maps $f_k\colon X_k\to Y_k$ is a map $F\colon X\to Y$ with monomorphisms
$\sigma_k\colon X_k\hookrightarrow X$, $\tau_k\colon Y_k\hookrightarrow Y$,
$F\sigma_k=\tau_kf_k$, such that $(\sigma_k\times\sigma_k)^{-1}(\ker F)=
\ker f_k$ and $\tau_k^{-1}(\im F)=\im f_k$ for every $k$, and such that
$F$ is $\bar c$-bidense whenever every $f_k$ is.
\end{defi}

\begin{thm}\label{thm:Bax}
Assume (H1), that every subobject has a witness, and that every sequence
of witnesses has a sum. Then $W_c\vee W_{c'}=W_{\bar c}$ iff $c^{(k)}=\bar c$
for some $k$.
\end{thm}
\begin{proof}
If $c^{(k)}=\bar c$, a $\bar c$-bidense map lies in $U_k\subseteq V$, and
$V\subseteq W_{\bar c}$ always. If $c^{(k)}\ne\bar c$ for every $k$, choose
for each $k$ a subobject $S_k\subseteq Y_k$ with $c^{(k)}(S_k)\ne\bar c(S_k)$;
replacing $Y_k$ by $\bar c(S_k)$ we may assume that $S_k$ is $\bar c$-dense in
$Y_k$ but $c^{(k)}(S_k)\ne Y_k$, as $c^{(k)}\le\bar c$ and both commute with
pullback along $\bar c(S_k)\hookrightarrow Y_k$. Let $f_k$ be a witness for
$S_k$ and $F$ a sum of the $f_k$. Each $f_k$ is $\bar c$-bidense, being a
regular epimorphism with $\ker f_k\subseteq\bar c(\Delta)$, so $F\in W_{\bar c}$.
But $F\notin U_k$ for every $k$: pulling back along $\sigma_k$ and $\tau_k$,
the conditions $c^{(k)}(\im F)=Y$ and $\ker F\subseteq c^{(k)}(\Delta_X)$
restrict, by pullback stability of $c^{(k)}$ and the definition of a sum,
to $c^{(k)}(\im f_k)=Y_k$ and $\ker f_k\subseteq c^{(k)}(\Delta_{X_k})$,
i.e.\ to $f_k\in U_k$, which fails since $c^{(k)}(S_k)\ne Y_k$. So
$F\in W_{\bar c}\setminus U=W_{\bar c}\setminus V$ by Theorem~\ref{thm:A}.
\end{proof}

It remains to produce witnesses and sums. The witness is always the
\emph{fold map} $f_S\colon Y\amalg_SY\to Y$ of the pushout of
$Y\hookleftarrow S\hookrightarrow Y$; what has to be checked is the description
of its kernel pair, and this is where the two classes of categories differ.

\subsection*{Coherent categories with effective unions.}
Assume
\begin{enumerate}
\item[(T)] finite unions of subobjects exist, are stable under pullback and
are \emph{effective}: $A\cup B$ is the pushout of $A\hookleftarrow A\cap B
\hookrightarrow B$; pushouts of monomorphisms along monomorphisms exist, are
pullback squares and have monic injections; and countable coproducts exist
and are disjoint and universal.
\end{enumerate}
Every Grothendieck topos satisfies (T) \cite[A1.4, A2.4, C2.2]{Elephant},
and (H1) as well, subobject lattices being frames. The category of groups
does not: pushouts of monomorphisms are amalgamated free products, which are
pullbacks with monic injections, but the union of the two factors in
$\mathbb Z*\mathbb Z$ is not the whole group, and coproducts are not
universal.

\begin{prop}\label{prop:witnessT}
Assume (T). For $S\subseteq Y$ the fold map $f_S\colon X_S:=Y\amalg_SY\to Y$
is a witness, and the coproduct $\coprod_kf_k$ of a sequence of witnesses
is a sum.
\end{prop}
\begin{proof}
Let $\iota_1,\iota_2\colon Y\hookrightarrow X_S$ be the injections, monomorphisms
with intersection $S$, which is also the equaliser of $\iota_1,\iota_2$.
The union $\iota_1Y\cup\iota_2Y\subseteq X_S$ is the pushout of
$\iota_1Y\hookleftarrow S\hookrightarrow\iota_2Y$ by effectiveness, hence the
inclusion $\iota_1Y\cup\iota_2Y\hookrightarrow X_S$ is an isomorphism: the two
injections \emph{cover} $X_S$. The map $f_S$ is a regular epimorphism,
being split by $\iota_1$. Since unions are pullback-stable, subobject
lattices are distributive and $X_S\times X_S=\bigcup_{a,b}\iota_aY\times
\iota_bY$; hence $\ker f_S=\bigcup_{a,b}\ker f_S\cap(\iota_aY\times\iota_bY)$,
and $\ker f_S\cap(\iota_aY\times\iota_bY)=(\iota_a\times\iota_b)
\bigl((\iota_a\times\iota_b)^{-1}\ker f_S\bigr)=(\iota_a\times\iota_b)(\Delta_Y)$,
as $f_S\iota_a=f_S\iota_b=\mathrm{id}_Y$. For $a=b$ this is contained in
$\Delta_{X_S}$; for $a\ne b$ it is the image of $(\iota_1,\iota_2)$ or
$(\iota_2,\iota_1)$. Now $(\iota_1,\iota_2)^{-1}(\Delta_{X_S})=S$, so by
pullback stability $(\iota_1,\iota_2)^{-1}(w(\Delta_{X_S}))=w(S)$, and the
image of $(\iota_1,\iota_2)$ lies in $w(\Delta_{X_S})$ iff $w(S)=Y$; the
same holds for $(\iota_2,\iota_1)$ by symmetry. So $\ker f_S\subseteq
w(\Delta_{X_S})$ iff $w(S)=Y$.

For the sum, let $F=\coprod_kf_k\colon\coprod X_k\to\coprod Y_k$ with the
coproduct injections $\sigma_k,\tau_k$, monomorphisms since coproducts are
disjoint. Since coproducts are universal, $(\coprod X_k)^2=\coprod_{k,l}
X_k\times X_l$, the pullback of $\ker F$ along $\sigma_k\times\sigma_k$ is
$\ker f_k$ and the pullback of $\im F$ along $\tau_k$ is $\im f_k$ (images
are stable under pullback). If every $f_k$ is $\bar c$-bidense then $F$ is a
regular epimorphism, and $\ker F$ is the union of the images of the
$\ker f_k$ (countable unions being images of coproducts), which lie in the
images of $\bar c(\Delta_{X_k})$, hence in $\bar c(\Delta_{\coprod X_k})$ by
Lemma~\ref{lem:calc}(i), as $\Delta_{\coprod X_k}$ is the union of the images
of the $\Delta_{X_k}$.
\end{proof}

\subsection*{Abelian categories.}
Let $\mathcal A$ be a Grothendieck abelian category; it is exact, satisfies
(H1) (subobject lattices are sets, unions of chains are directed colimits,
exact and pullback-stable by AB5, and distribute over finite intersections),
and has countable coproducts, but these are not disjoint and pushouts of
monomorphisms are not unions in the sense of (T). For a map $f\colon X\to Y$
write $d\colon X\times X\to X$ for $(x,x')\mapsto x-x'$, a split epimorphism
with $d^{-1}(0)=\Delta_X$ and $d^{-1}(K)=\ker f$ for the kernel object
$K=\operatorname{Ker}f\subseteq X$. For a word $w$, pullback stability gives
$w(\Delta_X)=d^{-1}(w(0))$, and since $d$ is a split epimorphism,
$d^{-1}(K)\subseteq d^{-1}(w(0))$ iff $K\subseteq w(0)$; so
\begin{equation}\label{eq:ab}
\ker f\subseteq w(\Delta_X)\quad\text{iff}\quad\operatorname{Ker}f\subseteq
w(0\subseteq X)\quad\text{iff}\quad w(0\subseteq\operatorname{Ker}f)=
\operatorname{Ker}f,
\end{equation}
the last step by pullback stability along $\operatorname{Ker}f\hookrightarrow
X$. Call an object $M$ \emph{$w$-torsion} if $w(0\subseteq M)=M$; by
pullback stability along a quotient map, $w(S)=Y$ iff $Y/S$ is $w$-torsion.

\begin{prop}\label{prop:witnessA}
In a Grothendieck abelian category, for $S\subseteq Y$ the fold map
$f_S\colon Y\amalg_SY\to Y$ is a witness, and the direct sum $\bigoplus_kf_k$
of a sequence of witnesses is a sum.
\end{prop}
\begin{proof}
$X_S=(Y\oplus Y)/\{(s,-s):s\in S\}$ and $f_S[(y_1,y_2)]=y_1+y_2$ is an
epimorphism with kernel object $\{[(y,-y)]\}\cong Y/S$. By \eqref{eq:ab},
$\ker f_S\subseteq w(\Delta_{X_S})$ iff $Y/S$ is $w$-torsion iff $w(S)=Y$. For
the sum, $F=\bigoplus f_k$ has $\operatorname{Ker}F=\bigoplus\operatorname{Ker}
f_k$ and $\operatorname{Coker}F=\bigoplus\operatorname{Coker}f_k$; the
pullbacks of $\ker F$ and $\im F$ along the summand inclusions are
$\ker f_k$ and $\im f_k$ by \eqref{eq:ab} and the corresponding statement
for images. If every $f_k$ is $\bar c$-bidense, then $F$ is an epimorphism
and $\operatorname{Ker}F=\bigoplus\operatorname{Ker}f_k$ is $\bar c$-torsion,
because $\bar c(0\subseteq\bigoplus K_k)\supseteq\bigoplus\bar c(0\subseteq
K_k)$ by Lemma~\ref{lem:calc}(i) applied to the summand inclusions and the
fact that unions of subobjects are sums; so $F\in W_{\bar c}$ by
\eqref{eq:ab}.
\end{proof}

\begin{proof}[Proof of Theorem~\ref{thm:B}]
Theorem~\ref{thm:Bax} with Lemma~\ref{lem:joinK} and
Proposition~\ref{prop:witnessT}, respectively~\ref{prop:witnessA}.
\end{proof}

\begin{figure}[ht]
\[
\begin{tikzcd}[column sep=large]
S\arrow[r,hook]\arrow[d,hook] & Y\arrow[d,"\iota_2"]\arrow[ddr,bend left=15,"\mathrm{id}"]\\
Y\arrow[r,"\iota_1"]\arrow[drr,bend right=15,"\mathrm{id}"'] & X_S\arrow[dr,"f_S"]\\
& & Y
\end{tikzcd}
\]
\caption{The witness: two copies of $Y$ glued along $S$; the kernel pair
of the fold map lies in $w(\Delta)$ exactly when $w(S)=Y$. In a topos the
two copies cover $X_S$; in an abelian category the kernel of the fold map is
$Y/S$.}
\end{figure}

\begin{rem}\label{rem:elem}
Theorem~\ref{thm:A} needs only exactness, and Proposition~\ref{prop:mono}
only finite limits. In Theorem~\ref{thm:Bax}, (H1) serves to identify
$\bar c$ with the transfinite iterate, witnesses to detect a single failure
and sums to combine infinitely many; without sums one still has, for each
fixed $k$, that $W_{\bar c}\subseteq U_k$ iff $c^{(k)}=\bar c$, provided
$\bar c$ exists and witnesses do.
\end{rem}

\section{Instances}\label{sec:instances}

\subsection{Toposes}\label{sub:topos}
Let $\cE$ be a Grothendieck topos. Closure operators on $\cE$ are the
universal closure operators, in bijection with Lawvere--Tierney topologies
$j\colon\Omega\to\Omega$ (the closure of $A\subseteq X$ is classified by
$j\circ\chi_A$) and with subtoposes, and $W_{c_j}$ is the class of maps
inverted by the sheafification $a_j$ \cite[V.1]{MLM}, \cite[A4.3]{Elephant}.
The hypotheses (H1) and (T) hold, and the identity $c^{(k)}=\bar c$ on all
subobject lattices is the identity $(jj')^k=j\vee j'$ on $\Omega$ (apply
both sides to the generic subobject). Theorems~\ref{thm:A} and~\ref{thm:B}
specialise to \cite[Thms.~1.1, 1.2]{Stabjoin}:

\begin{cor}\label{cor:topos}
For Lawvere--Tierney topologies $j,j'$ on a Grothendieck topos,
$W_j\vee W_{j'}$ is the class of maps whose image and kernel pair are reached
by a bounded number of alternating closures, the closure of $W_j\cup W_{j'}$
under composition, and $W_j\vee W_{j'}=W_{j\vee j'}$ iff $(jj')^k=j\vee j'$
for some $k$.
\end{cor}

For a presheaf topos $\Presh(S)$ and Grothendieck topologies $J,J'$ on the
small category $S$ the conditions are conditions on sieves: a map $f$ is in
$U_k$ iff its image sieves $\{b:y\cdot b\in\im f\}$ and its equaliser sieves
$\{b:x\cdot b=x'\cdot b\}$, $f(x)=f(x')$, are sent to the maximal sieve by
$(jj')^k$. For topologies of ideal type, $J_D$ and $J_E$ for idempotent
two-sided ideals $D,E$ of $S$, this is local bijectivity along the ideal
$(DE)^k$, and the criterion is the stabilisation of the chain $DE\supseteq
DEDE\supseteq\cdots$, which holds on every finite category and fails on
$(\mathbb N,\le)$ for the ideals of the even and of the odd points
\cite[\S5]{Stabjoin}. For ideals generated by two sets of objects $A,B$,
the join of ``bijective at the objects of $A$'' and ``bijective at the
objects of $B$'' is ``bijective at $A\cap B$'' iff every morphism through
a long enough alternating zigzag of $B$ and $A$ factors through $A\cap B$;
on $\mathbb N$ with the even and the odd points the join is the class of maps
of towers bijective up to a uniform shift, a proper subclass of the
isomorphisms of pro-sets \cite[\S6]{Stabjoin}.

\begin{rem}\label{rem:tower}
The classical sufficient condition for $W_j\vee W_{j'}=W_{j\vee j'}$ is that
the alternating tower of sheafifications $(a_{j'}a_j)^kX$ reach the
$(j\vee j')$-sheaves for one $k$ and all $X$ \cite[Thm.~4.1]{Postnikov}. The
criterion of Theorem~\ref{thm:B} is the iteration of the closure operators,
and the two differ: on the contracted monoid of order $8$ of
\cite[\S7]{Finitejoin} the chain stabilises, so the join is preserved, while
the strictly alternating plus tower never reaches the sheaves uniformly
\cite[Prop.~7.1]{Finitejoin}, and a computation on the right $M$-sets with
at most three elements shows that no composite of at most eight plus
constructions, in particular no $(a_Ea_D)^k$ with $k\le2$, is a sheaf for all
of them \cite[\S8]{Stabjoin}.
\end{rem}

\subsection{Locales}\label{sub:locale}
Let $L$ be a frame and $\Sh(L)$ its topos of sheaves. Closure operators on
$\Sh(L)$ correspond to the nuclei on $L$ (sublocales), ordered as the
nuclei \cite[II.2]{StoneSpaces}, \cite[A4.5]{Elephant}; the subobjects of
the representable $y(U)$ are the $V\le U$; and for the nucleus $n$ with
sublocale inclusion $i$, the sheafification is $i_*i^*$ with $i^*(y(V))=
y(n(V))$, so $V\hookrightarrow W$ is dense iff $n(V)=n(W)$, and the closure of
$V$ in $y(U)$ is $n(V)\wedge U$, the largest $W\le U$ with $n(W)=n(V)$. By
induction, using $U\le n(U),n'(U)$, $(cc')^k(V)=(nn')^k(V)\wedge U$ and
$\bar c(V)=(n\vee n')(V)\wedge U$ \cite[Lemma~7.1]{Stabjoin}.

\begin{cor}\label{cor:localic}
For nuclei $n,n'$ on a frame $L$, $W_n\vee W_{n'}=W_{n\vee n'}$ in
$\Same(\Sh(L))$ iff $(nn')^k=n\vee n'$ in the frame of nuclei for some $k$.
\end{cor}
\begin{proof}
Theorem~\ref{thm:B} and the formula above with $U=\top$.
\end{proof}

Open and closed nuclei satisfy the condition with $k=1$: $o(U)c(U')$,
$V\mapsto U\to(U'\vee V)$, is idempotent and hence the join. For a
topological space $X$ and subspaces $A,B$, with nuclei $n_A(V)=
\operatorname{int}(V\cup(X\setminus A))$, the class $W_{n_A}$ consists of the
maps of sheaves whose inverse image to $A$ is an isomorphism, and the
condition reads: the alternating derivation $T(F)=\operatorname{cl}
(\operatorname{cl}(F\cap B)\cap A)$ on closed sets reaches its transfinite
limit in a uniform number of steps \cite[Prop.~7.5]{Stabjoin}. It holds for
closed and for open subspaces with $k=1$ and fails for the rationals and
the irrationals in $\mathbb R$: for a countable closed set $C$ whose points
of even Cantor--Bendixson rank are rational and of odd rank irrational,
$T^j(C)$ is the $2j$-th derived set, nonempty for $2j$ below the rank, while
the transfinite limit is empty; so in $\Same(\Sh(\mathbb R))$ the join of
``isomorphism over $\mathbb Q$'' and ``isomorphism over the irrationals'' is
strictly smaller than the class of the meet subtopos, although the two
sublocales meet in a nonempty dense sublocale \cite[Ex.~7.6]{Stabjoin}.

\subsection{Frames as categories: the pointwise criterion}\label{sub:frame}
Let $L$ be a frame, or just a meet-semilattice with top, viewed as a
category. It is exact: every morphism $x\le y$ is a monomorphism, the only
regular epimorphisms are the identities, and the (regular epi, mono)
factorisation of $x\le y$ is $x\le y$ itself; kernel pairs are diagonals and
equivalence relations trivial. $\Sub(y)=\mathord\downarrow y$, and the
pullback of $z\le y$ along $x\le y$ is $z\wedge x$. For a nucleus $n$ on $L$
put $c_n(z\le y):=n(z)\wedge y$; this is a closure operator: inflationary,
monotone, idempotent ($n(n(z)\wedge y)\wedge y=n(z)\wedge n(y)\wedge y=n(z)
\wedge y$), intersection-preserving, and pullback-stable
($c_n(z\wedge x)=n(z)\wedge n(x)\wedge x=c_n(z)\wedge x$). The $c_n$-bidense
maps are the $x\le y$ with $n(x)\wedge y=y$, i.e.\ $n(x)\ge y$, i.e.\
$n(x)=n(y)$: $W_{c_n}$ is the class $W_n^L=\{x\le y:nx=ny\}$ of
\cite[Thm.~1.1]{Fuzzysame}.

\begin{cor}\label{cor:frame}
For nuclei $n,n'$ on $L$, in $\Same(L)$,
\[
W_n^L\vee W_{n'}^L=\{x\le y:\ (nn')^kx\ge y\text{ for some }k\},
\]
every such $x\le y$ being the composite of the chain $x\le n'x\wedge y\le
nn'x\wedge y\le\cdots\le y$ of inequalities alternately in $W_{n'}^L$ and
$W_n^L$; and $W_n^L\vee W_{n'}^L=W_{n\vee n'}^L$ iff for every $x$ some
iterate $(nn')^kx$ equals $(n\vee n')x$.
\end{cor}
\begin{proof}
Theorem~\ref{thm:A}: $x\le y$ is in $U_k$ iff $c^{(k)}(x)=(nn')^kx\wedge y=y$,
and the peeling factorisation is the displayed chain. For the last
assertion, $W^L_{n\vee n'}=\{x\le y:(n\vee n')x\ge y\}$; if the pointwise
identity holds, then $(n\vee n')x\ge y$ gives $(nn')^kx\ge y$ for the $k$ of
$x$; conversely take $y=(n\vee n')x$.
\end{proof}

Here Theorem~\ref{thm:Bax} does not apply: in a poset the pushout of
$y\hookleftarrow s\hookrightarrow y$ is $y$ with both injections the identity,
its fold map is $\mathrm{id}_y$, and no witnesses exist (the only regular
epimorphisms are identities, whose kernel pair is the diagonal). Accordingly the criterion is pointwise, as in
\cite[Thm.~1.1]{Fuzzysame}, and the examples there of nuclei on the
G\"odel chain $[0,1]$ that stabilise pointwise but not uniformly separate
$\Same(L)$ from $\Same(\Sh(L))$: with $F=\{0,1\}\cup\{1/2k\}$ and
$G=\{0,1\}\cup\{1/(2k+1)\}$ as fixed sets, $W_{n_F}\vee W_{n_G}=W_{n_F\vee
n_G}$ in $\Same(L)$ by Corollary~\ref{cor:frame} but not in
$\Same(\Sh([0,1]))$ by Corollary~\ref{cor:localic} \cite[Ex.~7.4]{Stabjoin}.

\subsection{Fuzzy sets}\label{sub:fuzzy}
Let $L$ be a complete lattice and $\Set(L)$ Goguen's category of $L$-fuzzy
sets: objects $(X,\mu)$ with $\mu\colon X\to L$, morphisms $f\colon(X,\mu)
\to(Y,\nu)$ maps with $\mu\le\nu f$. It has finite limits (computed on
underlying sets with the initial structure) but is not exact: for the
three-element chain $L=\{0<a<1\}$ and $X=\{x_0,x_1\}$ with $\mu\equiv1$, the
relation on $X\times X$ with degree $1$ on the diagonal and $a$ off it is an
internal equivalence relation that is the kernel pair of no map, as kernel
pairs of maps out of $(X,\mu)$ have degree $1$ wherever they are nonempty.
For a nucleus $j$ on $L$ preserving binary meets, the $j$-closed fuzzy sets
($\mu=j\mu$) form a reflective subcategory with reflector $(X,\mu)\mapsto
(X,j\mu)$, and the class $W_j$ of maps inverted by it consists of the
bijections $f$ with $j\mu=j\nu f$ \cite{Fuzzysame}. Subobjects of $(Y,\nu)$
are the $(A,\alpha)$ with $A\subseteq Y$ and $\alpha\le\nu|_A$, pulled back
along $f$ to $(f^{-1}A,\alpha f\wedge\mu)$; put
\[
c_j(A,\alpha):=(Y,\ j\bar\alpha\wedge\nu),
\]
where $\bar\alpha$ extends $\alpha$ by the bottom element outside $A$. This
is inflationary, monotone, idempotent, preserves binary intersections (as
$j$ preserves binary meets) and is pullback-stable: both
$f^{-1}(c_j(A,\alpha))$ and $c_j(f^{-1}(A,\alpha))$ are
$(X,j(\bar\alpha f)\wedge\mu)$, using $\mu\le\nu f$ and $\mu\le j\mu$. A
bijection $f\colon(X,\mu)\to(Y,\nu)$ is a monomorphism with image
$(Y,\mu f^{-1})$, so it is $c_j$-dense iff $j(\mu f^{-1})\ge\nu$ iff
$j\mu=j\nu f$: on bijections, $c_j$-density is membership in $W_j$, and
$f\in M_k$ iff $(jj')^k\mu\ge\nu f$. Bijections form a sameness $B$, and for
a bijection $f$ the inclusions $\im f\hookrightarrow c_j(\im f)$ are
bijections, since $c_j(\im f)$ has underlying set $Y$. Hence
Proposition~\ref{prop:mono} gives
\[
W_j\vee W_{j'}=\{f\text{ bijective}:\ (jj')^k\mu\ge\nu f\text{ for some }k\},
\]
the description of \cite[Lemma~4.1, Thm.~4.3]{Fuzzysame} by uniformly
bounded alternating chains, each such $f$ being a composite of bijections
alternately in $W_{j'}$ and $W_j$. The uniform criterion
$(j'j)^k=j\vee j'$ of \cite[Thm.~1.1(ii)]{Fuzzysame} is proved there with
witnesses that glue fuzzy sets along a common part and sums of such
witnesses, in the pattern of Theorem~\ref{thm:Bax}; we do not claim
$\Set(L)$ as an instance of Theorem~\ref{thm:B}. (For non-bijective maps the
class of $c_j$-bidense maps is larger than $W_j$: the fold map of two points
of degree $\bot$ onto one is $c_j$-bidense for every $j$. The join of two
classes of bijections is computed among bijections, so this does not affect
the statement.)

\subsection{Abelian categories and torsion theories}\label{sub:abelian}
Let $\mathcal A$ be a Grothendieck abelian category. A \emph{hereditary
torsion theory} is given by a class $T$ of objects closed under subobjects,
quotients, extensions and coproducts; equivalently by its radical
$t\colon\mathcal A\to\mathcal A$, the subfunctor of the identity with $t(M)$ the
largest subobject of $M$ in $T$, which is left exact ($t(N)=N\cap t(M)$ for
$N\subseteq M$) and idempotent ($t(M/t(M))=0$) \cite{Gabriel,Stenstrom}.

\begin{prop}\label{prop:torsion}
Closure operators on $\mathcal A$ correspond bijectively to hereditary
torsion theories (this is classical, cf.\ \cite{DT,Stenstrom}; we include
the proof to fix the dictionary): $T$ gives $c_T(A):=\pi^{-1}(t(X/A))$ for $A\subseteq X$ with
$\pi\colon X\to X/A$, and $c$ gives $T_c=\{M:c(0\subseteq M)=M\}$ with radical
$t_c(M)=c(0\subseteq M)$. Under this correspondence $c_T\le c_{T'}$ iff
$T\subseteq T'$, $W_{c_T}$ is the class of maps with kernel and cokernel in
$T$, the maps inverted by the quotient functor $\mathcal A\to\mathcal A/T$,
and the least closure operator above $c_T$ and $c_{T'}$ is $c_{T\vee T'}$ for
the least hereditary torsion theory $T\vee T'$ containing $T$ and $T'$.
\end{prop}
\begin{proof}
$c_T$ is inflationary, monotone (the image of $t(X/A)$ in $X/B$ lies in
$t(X/B)$ for $A\subseteq B$) and idempotent ($X/c_T(A)=(X/A)/t(X/A)$ is
torsion-free). It preserves intersections: $X/(A\cap B)$ embeds in
$X/A\oplus X/B$, and as $t$ is left exact and additive, $x+(A\cap B)$ is
torsion iff $x+A$ and $x+B$ are. It is pullback-stable: for $f\colon X'\to X$,
$X'/f^{-1}(A)$ embeds in $X/A$, so $t(X'/f^{-1}(A))$ is the preimage of
$t(X/A)$ by left exactness, i.e.\ $c_T(f^{-1}A)=f^{-1}(c_TA)$. Conversely, for
a closure operator $c$, $t_c(M)=c(0)$ is a subfunctor (by
Lemma~\ref{lem:calc}(i)), left exact (pullback stability along
monomorphisms), and idempotent ($c(0\subseteq M/c(0))$ pulls back to $c(c(0))=
c(0)$), so $T_c$ is a hereditary torsion theory, and $c(A)=\pi^{-1}(c(0
\subseteq X/A))=c_{T_c}(A)$ by pullback stability along $\pi$; and
$t_{c_T}=t$. Order: $c_T\le c_{T'}$ iff $t\le t'$ iff $T\subseteq T'$. A map
$f\colon X\to Y$ has $c_T(\im f)=Y$ iff $Y/\im f=\operatorname{Coker}f$ is
torsion, and $\ker f\subseteq c_T(\Delta_X)$ iff $\operatorname{Ker}f\subseteq
t(X)$ iff $\operatorname{Ker}f\in T$ by \eqref{eq:ab}; these are the maps
inverted by the quotient functor \cite[III]{Gabriel}. The last assertion
follows from the order isomorphism.
\end{proof}

For a word $w$ in $c_T,c_{T'}$ the $w$-torsion objects are those with a
filtration by the canonical subobjects $0\subseteq t'(M)\subseteq
\pi^{-1}t(M/t'(M))\subseteq\cdots$; so an object is $c^{(k)}$-torsion iff it
lies in the class $E_k$ of Corollary~\ref{cor:abelian}, any alternating
filtration being dominated by the canonical one, and $E_k$ is closed under
subobjects and quotients. By \eqref{eq:ab} and pullback stability along
$Y\to Y/\im f$, a map is in $U_k$ iff its kernel and cokernel lie in $E_k$.

\begin{proof}[Proof of Corollary~\ref{cor:abelian}]
Theorems~\ref{thm:A} and~\ref{thm:B} with Proposition~\ref{prop:torsion}
and the description of $U_k$; the identity $c^{(k)}=\bar c$ says that every
$c_{T\vee T'}$-torsion object, i.e.\ every object of $T\vee T'$, is
$c^{(k)}$-torsion.
\end{proof}

By Lemma~\ref{lem:joinK}, $T\vee T'$ is the class of objects with a
transfinite alternating filtration, the classical description of the join
of two hereditary torsion theories; the corollary says that the join of
the two classes of maps inverted by the Gabriel quotients is the class of
the join torsion theory exactly when a finite bound on the length exists.
When $T$ and $T'$ commute in the sense that $t(t'(M))=t'(t(M))$ and the
join is the class of extensions of a $T$-object by a $T'$-object, $k=1$:
for instance the $p$-torsion and the $q$-torsion abelian groups for distinct
primes.

\begin{ex}\label{ex:persistence}
Let $\mathcal A$ be the category of persistence modules over a field $k$:
representations $\cdots\to M_2\to M_1\to M_0$ of the ordinal $\omega$ with
arrows $n+1\to n$, a Grothendieck abelian category. Let $T$ and $T'$ be the
classes of modules supported on the even, respectively on the odd, vertices;
each is closed under subobjects, quotients, extensions and direct sums,
hence a hereditary torsion theory, with $t(M)_n=0$ for $n$ odd and
$t(M)_n=\operatorname{Ker}(M_n\to M_{n-1})$ for $n$ even ($t(M)_0=M_0$), and
symmetrically $t'$. A map $f$ has kernel and cokernel in $T$ iff $f_n$ is an
isomorphism for every odd $n$, so $W_T$ is the class of maps that are
isomorphisms at the odd vertices and $W_{T'}$ at the even vertices, the
linear analogue of the classes of \cite[Thm.~7.1]{Topjoin}. The join
torsion theory is everything: a module that is torsion-free for both has
$M_0=0$ and all maps $M_n\to M_{n-1}$ injective, hence is $0$. The
interval module $I=k[0,N]$, with $I_n=k$ for $n\le N$, identities as maps
and $I_n=0$ for $n>N$, has $t'(I)=0$, $t(I)=k[0,0]$, then $t'(I/k[0,0])=
k[1,1]$, and so on: its canonical alternating filtration is
$0\subseteq0\subseteq k[0,0]\subseteq k[0,1]\subseteq\cdots\subseteq k[0,N]$,
of length $N+2$, and $I\in E_k$ iff $2k\ge N+2$. So $E_k\ne T\vee T'$ for
every $k$, and by Corollary~\ref{cor:abelian} the join of ``isomorphism at
the odd vertices'' and ``isomorphism at the even vertices'' is a proper
class of maps, although the join localisation is the zero category and
$W_{T\vee T'}$ is the class of all maps; it is the class of maps whose kernel
and cokernel have alternating filtrations of bounded length, for instance
it contains $0\to k[0,N]$ for each $N$ but not $0\to\bigoplus_Nk[0,N]$. The
witness of Proposition~\ref{prop:witnessA} for $S=0\subseteq k[0,N]$ is the
fold map $k[0,N]\oplus k[0,N]\to k[0,N]$, in $W_{T\vee T'}$ but not in
$U_k$ for $2k<N+2$, and the direct sum of these fold maps over $N$ lies in
no $U_k$.
\end{ex}

\subsection{Summary}
The settings are instances of one theorem, and the identity
$(cc')^k=\bar c$ that governs them is, for nuclei, the finite-iteration
axiom of the lax logic of the join of two nuclei \cite{Laxjoin}: uniform in
the subobject in toposes, locales, abelian categories and fuzzy sets, where
witnesses and sums exist, and pointwise in the frame itself, where they do
not. The pointwise condition is also the one that characterises
$\Box_C\vee\Box_{C'}=\Box_{C\cap C'}$ for relativized necessities on the
frame of up-sets of a poset (absence of infinite alternating chains
\cite{Relnec}), and a disjoint union of finite alternating chains of
unbounded length satisfies it but not the uniform one.

\section{The dual regime: quotient-closed reflections}\label{sec:dual}

Let $\mathcal V$ be a variety of algebras and $R\subseteq\mathcal V$ a
subvariety, with reflector $a_R(A)=A/\rho_R(A)$, $\rho_R(A)$ the least
congruence with quotient in $R$. A homomorphism $f\colon A\to B$ is
inverted by $a_R$ iff $\im f$ meets every $\rho_R(B)$-class and
$(f\times f)^{-1}(\rho_R(B))\subseteq\rho_R(A)$. This looks like bidensity,
but the candidate closure operators are not pullback-stable: for
$R$ the abelian groups, $\rho_R(A)$ is the commutator congruence, and the
inclusion $A_3\hookrightarrow S_3$ has $\rho_R(A_3)$ trivial while
$\rho_R(S_3)$ restricted to $A_3$ is everything; likewise the
$\rho_R(B)$-saturation of a subalgebra is not stable under intersection
with subalgebras (in $S_3$, with $K=\{e,(13)\}$ and $A=\{e,(12)\}$, the
saturation of $A$ contains $(13)$ but $A\cap K$ is trivial). Accordingly
the mechanism is different: because $R$ is closed under quotients, the
composite $a_{R'}a_R$ is already the reflector onto $R\cap R'$, and the
join theorem of \cite[Thm.~4.1]{Postnikov} gives $W_R\vee W_{R'}=
W_{R\cap R'}$ for all subvarieties, with no condition
\cite[Thm.~1.1]{Varieties}: at the level of reflectors, $k=1$. The hard
direction there is the meet, $W_R\cap W_{R'}=W_{R\vee R'}$, which holds iff
a fracture condition is satisfied (subdirect representation and absence of
dense proper subalgebras \cite[Thm.~1.2]{Varieties}). The two regimes are
thus dual: for closure-type reflections (subtoposes) the meet of classes
is automatic and the join is governed by finite iteration of closures;
for quotient-closed reflections (subvarieties) the join is automatic and
the meet is governed by fracture. Remark~\ref{rem:tower} shows that in the
closure-type regime the reflector-level condition ($k=1$ or any uniform
$k$ for the tower of reflectors) is not the right one.

\section{Questions}

\begin{que}
Theorem~\ref{thm:B} reduces join preservation to the identity
$(cc')^k=\bar c$. For subspaces this is the uniform termination of the
alternating derivation on closed sets, for ideal topologies the
stabilisation of the ideal chain. Is there a structural description of the
pairs of nuclei, or of subtoposes, satisfying it, for instance in terms of
the lattice of sublocales?
\end{que}

\begin{que}
Does Theorem~\ref{thm:B} hold in an elementary topos, where (H1) and the
coproducts in (T) may be unavailable? By Remark~\ref{rem:elem} the obstacle
is the existence of $\bar c$ and the combination of infinitely many
witnesses. Likewise, which exact categories beyond toposes and abelian
categories have witnesses and sums?
\end{que}

\begin{que}
Is there a common generalisation of the closure-type and the
quotient-closed regimes, for reflective subcategories of an exact category
whose reflector is determined by a closure operator that is only
\emph{weakly} hereditary in the sense of \cite{DT}?
\end{que}

\section*{Funding}
No funding was received for this work.

\section*{Declarations}
The author has no competing interests to declare.

\begin{thebibliography}{9}
\bibitem{Barr} M.~Barr, \emph{Exact categories}, in: Exact Categories and
Categories of Sheaves, Lecture Notes in Math. 236, Springer, 1971, 1--120.
\bibitem{CDT} M.~M.~Clementino, D.~Dikranjan, W.~Tholen, \emph{Torsion
theories and radicals in normal categories}, J. Algebra \textbf{305} (2006),
98--129.
\bibitem{DT} D.~Dikranjan, W.~Tholen, \emph{Categorical Structure of Closure
Operators}, Mathematics and its Applications 346, Kluwer, 1995.
\bibitem{Elephant} P.~T.~Johnstone, \emph{Sketches of an Elephant: A Topos
Theory Compendium}, Oxford Logic Guides 43--44, Oxford University Press, 2002.
\bibitem{Gabriel} P.~Gabriel, \emph{Des cat\'egories ab\'eliennes}, Bull.
Soc. Math. France \textbf{90} (1962), 323--448.
\bibitem{Dickson} S.~E.~Dickson, \emph{A torsion theory for abelian
categories}, Trans. Amer. Math. Soc. \textbf{121} (1966), 223--235.
\bibitem{Escardo} M.~H.~Escard\'o, \emph{Joins in the frame of nuclei}, Appl.
Categ. Structures \textbf{11} (2003), 117--124.
\bibitem{StoneSpaces} P.~T.~Johnstone, \emph{Stone Spaces}, Cambridge
Studies in Advanced Mathematics 3, Cambridge University Press, 1982.
\bibitem{KM} Z.~Khanjanzadeh, A.~Madanshekaf, \emph{Weak topologies on
toposes}, Bull. Iranian Math. Soc. \textbf{47} (2021), 461--486;
arXiv:1808.06104.
\bibitem{MLM} S.~Mac Lane, I.~Moerdijk, \emph{Sheaves in Geometry and Logic},
Springer, 1992.
\bibitem{Same} R.~Shibusawa, \emph{The sameness lattice of a category:
weak-equivalence structures, Galois functoriality, and comparison intervals},
preprint (2026), archived at DOI 10.5281/zenodo.22823772.
\bibitem{Postnikov} R.~Shibusawa, \emph{Postnikov truncations as joins of sheaf
and homotopy equivalences}, preprint (2026), archived at DOI 10.5281/zenodo.22956181.
\bibitem{Topjoin} R.~Shibusawa, \emph{Meets and joins of Grothendieck
topologies in the sameness lattice}, preprint (2026), archived at DOI
10.5281/zenodo.22978446.
\bibitem{Finitejoin} R.~Shibusawa, \emph{Grothendieck topologies on a finite
category form a sublattice of the lattice of two-out-of-three classes},
preprint (2026), archived at DOI 10.5281/zenodo.23204943.
\bibitem{Varieties} R.~Shibusawa, \emph{Lattices of subvarieties and of
logics as sublattices of the lattice of two-out-of-three classes}, preprint
(2026), archived at DOI 10.5281/zenodo.23407364.
\bibitem{Laxjoin} R.~Shibusawa, \emph{The join of two nuclei as a lax modality:
axiomatization, finite models, and inexpressibility}, preprint (2026),
archived at DOI 10.5281/zenodo.23330576.
\bibitem{Relnec} R.~Shibusawa, \emph{Joins of relativized necessities on a
poset: an alternating-chain criterion and the induction scheme of lax logic},
preprint (2026), archived at DOI 10.5281/zenodo.23411203.
\bibitem{Fuzzysame} R.~Shibusawa, \emph{Nuclei on the algebra of truth values
and the sameness lattice of fuzzy sets: meets, joins and the
finite-iteration axiom}, preprint (2026), archived at DOI 10.5281/zenodo.23413073.
\bibitem{Stenstrom} B.~Stenstr\"om, \emph{Rings of Quotients}, Grundlehren
der mathematischen Wissenschaften 217, Springer, 1975.
\bibitem{Stabjoin} R.~Shibusawa, \emph{Joins of Lawvere--Tierney topologies
in the lattice of two-out-of-three classes: the finite-iteration criterion},
preprint (2026), archived at DOI 10.5281/zenodo.23487202.
\end{thebibliography}

\end{document}
